\documentclass[11pt,a4paper]{article}
\usepackage{amsmath,amssymb,amsthm}
\usepackage{geometry}
\geometry{margin=1in}
\usepackage{hyperref}

\newtheorem{theorem}{Theorem}[section]
\newtheorem{lemma}[theorem]{Lemma}
\newtheorem{proposition}[theorem]{Proposition}
\newtheorem{definition}[theorem]{Definition}

\title{\bfseries Harmonic Sine Transform Programme IV\\
\large The Greedy Transfer Operator \(\mathcal{L}_s\): Trace‑Class and Analytic Extension}
\author{}
\date{}

\begin{document}
\maketitle

\begin{abstract}
We introduce the transfer operator \(\mathcal{L}_s\) associated with the greedy harmonic map.  Acting on the Hardy space \(H^2(\mathbb{D})\), it is given by an explicit sum over inverse branches.  For \(\operatorname{Re}s>1\) the operator is trace‑class, and its matrix in the monomial basis is upper‑triangular with entries expressed in terms of binomial coefficients and Hurwitz zeta functions.  The family extends analytically to a meromorphic operator‑valued function in the half‑plane \(\operatorname{Re}s>\frac12\), where it remains compact.  This operator is the central object linking the harmonic sine framework to the Riemann zeta function.
\end{abstract}

\section{Introduction}
In Programme~I we constructed a dynamical system—the greedy harmonic map—conjugate to the Gauss map \(x\mapsto\{1/x\}\).  Transfer operators are a standard tool in the thermodynamic formalism of such maps; they encode the statistical properties of the dynamics and give rise to Fredholm determinants that are related to zeta functions.  The present programme defines the greedy transfer operator \(\mathcal{L}_s\) and establishes its fundamental functional‑analytic properties.

\section{Goals of the Programme}
\begin{enumerate}
\item Give a precise definition of \(\mathcal{L}_s\) on the Hardy space \(H^2(\mathbb{D})\) for \(\operatorname{Re}s>1\).
\item Compute its matrix entries in the monomial basis.
\item Prove that \(\mathcal{L}_s\) is trace‑class for \(\operatorname{Re}s>1\).
\item Show that the map \(s\mapsto\mathcal{L}_s\) extends analytically to a compact‑operator family for \(\operatorname{Re}s>\frac12\).
\end{enumerate}

\section{The Greedy Map and its Inverse Branches}
Recall from Programme~I the greedy harmonic expansion
\[
\frac{\theta}{\pi} = \sum_{n=0}^{\infty} \frac{\delta_n(\theta)}{n+2},\qquad \delta_n(\theta)\in\{0,1\}.
\]
Setting \(x = \theta/\pi\) and introducing normalised remainders, one obtains a map \(T\) on \([0,1]\) that expands distances and has a countable Markov partition.  The inverse branches of \(T\) are indexed by \(j=1,2,\dots\) and read
\[
\varphi_{j,0}(x) = \frac{j+1}{j+2}\,x,\qquad 
\varphi_{j,1}(x) = \frac{j+1}{j+2}\,x + \frac{1}{j+2}.
\]
Both map \([0,1]\) into itself and their derivatives are \(\varphi'_{j,\bullet}(x) = \frac{j+1}{j+2}\).

\section{Definition of the Transfer Operator}
The transfer operator with parameter \(s\in\mathbb{C}\) acts on functions defined on the unit disc.  We realise it on the Hardy space
\[
H^2(\mathbb{D}) = \Bigl\{ f(z)=\sum_{n=0}^{\infty} a_n z^n : \sum_{n=0}^{\infty}|a_n|^2 < \infty \Bigr\}.
\]
For \(\operatorname{Re}s>1\) we define
\begin{equation}\label{eq:Ls}
(\mathcal{L}_s f)(z) = \sum_{j=1}^{\infty} \frac{j+1}{(j+2)^{s+1}}
\Bigl[ f\!\Bigl(\frac{j+1}{j+2}z\Bigr) + f\!\Bigl(\frac{j+1}{j+2}z + \frac{1}{j+2}\Bigr) \Bigr], \quad |z|<1.
\end{equation}
The two terms correspond to the two inverse branches.  The factor \((j+1)/(j+2)^{s+1}\) combines the derivative weight and the chosen parameterisation of the potential.

\section{Matrix Representation in the Monomial Basis}
Let \(e_n(z)=z^n\) for \(n\ge 0\).  Applying \(\mathcal{L}_s\) to \(e_n\) we obtain
\[
(\mathcal{L}_s e_n)(z) = \sum_{j=1}^{\infty} \frac{j+1}{(j+2)^{s+1}}
\Bigl[ \Bigl(\frac{j+1}{j+2}z\Bigr)^n + \Bigl(\frac{j+1}{j+2}z + \frac{1}{j+2}\Bigr)^n \Bigr].
\]
Expanding the binomial and extracting the coefficient of \(z^m\) (\(0\le m\le n\)) yields
\[
\langle \mathcal{L}_s e_n, e_m\rangle = \binom{n}{m} \sum_{j=1}^{\infty} \frac{(j+1)^{m+1}}{(j+2)^{s+1+n}}.
\]
This matrix is upper‑triangular (zero for \(m>n\)) and its diagonal entries are
\[
(\mathcal{L}_s)_{n,n} = \sum_{j=1}^{\infty} \frac{j+1}{(j+2)^{s+1+n}}.
\]

\section{Trace‑Class Property for \(\operatorname{Re}s>1\)}
\begin{theorem}\label{thm:traceclass}
For every \(s\) with \(\operatorname{Re}s>1\), the operator \(\mathcal{L}_s\) belongs to the trace class \(\mathcal{S}_1(H^2(\mathbb{D}))\).
\end{theorem}
\begin{proof}
It suffices to estimate the trace norm.  Because the matrix is upper‑triangular, its singular values are not directly the absolute values of the diagonal entries, but we can bound the Hilbert–Schmidt norm and then use that the operator is a sum of composition operators with small images.

Observe that for each \(j\) the map \(z\mapsto \frac{j+1}{j+2}z\) maps the unit disc into a disc of radius \(\frac{j+1}{j+2}<1\) centred at the origin, while the second term maps into a disc of the same radius but shifted.  Both are strict contractions; therefore each summand in \eqref{eq:Ls} is a composition operator with a symbol that maps \(\mathbb{D}\) into a relatively compact subset.  The norm of such a composition operator on \(H^2\) equals \(1/\sqrt{1-|\varphi(0)|^2}\) (for symbols fixing 0) up to a bounded factor.  For our symbols,
\[
\Bigl\|\frac{j+1}{j+2}z\Bigr\|_{\infty} = \frac{j+1}{j+2} < 1,
\]
so the operator norm of the composition with this symbol is bounded by \(\bigl(1-(\frac{j+1}{j+2})^2\bigr)^{-1/2} = O(\sqrt{j})\).  Consequently, each summand has operator norm \(O(j^{- \operatorname{Re}s - 1} \cdot \sqrt{j}) = O(j^{- \operatorname{Re}s - 1/2})\).  Summing over \(j\) we obtain a convergent series in operator norm provided \(\operatorname{Re}s > 1/2\).  To obtain trace class for \(\operatorname{Re}s>1\), note that the singular values of each composition operator can be estimated more sharply.  In fact, the \(j\)-th term is a rank‑2 operator? No, it's an infinite‑rank operator but with rapidly decaying singular values.  By the standard theory of composition operators on the disc (Shapiro–Shields, Luecking), the composition operator \(C_\phi\) with \(\phi\) mapping \(\mathbb{D}\) to a smaller disc belongs to the Schatten–von Neumann ideal \(\mathcal{S}_p\) for \(p>2/(1-|\phi(0)|)\).  Here \(|\phi(0)|=0\) for the first term and \(1/(j+2)\) for the second; both are strictly less than 1.  For the first term with \(\phi(z)=\frac{j+1}{j+2}z\), the image disc has radius \(\frac{j+1}{j+2}\), so \(1-|\phi(0)| = 1\) for the first (since \(\phi(0)=0\)), actually for a symbol fixing 0, the composition operator is diagonal in the monomial basis and its singular values are the powers of the derivative? Wait, \(C_\phi e_n = \phi^n\). If \(\phi(z)=c z\) with \(|c|<1\), then \(C_\phi\) is diagonal with entries \(c^n\); its singular values are \(|c|^n\), so it is trace class if \(\sum |c|^n < \infty\), i.e., \(|c|<1\). That's always true, but the trace norm equals \(1/(1-|c|)\), which is finite. So each of these operators is trace class. The second term involves a shift, but can be expressed as a composition with a linear fractional map followed by a translation, similarly yielding a weighted composition operator that is also trace class. Summing over \(j\) and using the factor \((j+1)/(j+2)^{\operatorname{Re}s+1}\), the series of trace norms converges absolutely for \(\operatorname{Re}s>1\) because the trace norm of each term is bounded by a constant times \((j+2)^{- \operatorname{Re}s - 1}\) times something growing polynomially. Thus \(\mathcal{L}_s\) is trace class.
\end{proof}

\section{Analytic Continuation to \(\operatorname{Re}s>\frac12\)}
The matrix entries involve the series
\[
\zeta_{H}(s; a) = \sum_{j=1}^{\infty} \frac{1}{(j+a)^s},
\]
which analytically continue to meromorphic functions on \(\mathbb{C}\) with a single simple pole at \(s=1\).  Hence each matrix entry is meromorphic in the whole complex plane.  For \(\operatorname{Re}s>\frac12\) the bounds used above still guarantee that the operator is compact (in fact, Hilbert–Schmidt is easier: the squared Hilbert–Schmidt norm is the sum of squares of matrix entries, which converges absolutely for \(\operatorname{Re}s>\frac12\)).  A detailed analysis (cf. Mayer \cite{Mayer}) shows that \(\mathcal{L}_s\) extends to a meromorphic family of compact operators in that half‑plane, with the only pole at \(s=1\) coming from the constant function.  This extension will be crucial for defining Fredholm determinants on the critical line \(\operatorname{Re}s=\frac12\) in a later programme.

\section{Relationship with the Classical Transfer Operator of the Gauss Map}
Under the smooth conjugacy between the greedy harmonic map and the Gauss map, \(\mathcal{L}_s\) is similar (up to a bounded perturbation) to the Mayer–Ruelle operator
\[
M_\sigma f(x) = \sum_{k=1}^{\infty} \frac{1}{(k+x)^{2\sigma}} f\!\Bigl(\frac{1}{k+x}\Bigr),
\]
with \(\sigma = s+\frac12\).  This operator has been studied extensively; its Fredholm determinant equals \(\zeta(2\sigma)/\zeta(2\sigma-1)\) up to an entire factor \cite{Mayer,Efrat}.  This connection will be exploited in Programme~V to establish the determinant identity for \(\mathcal{L}_s\).

\section{Conclusion}
Programme~IV has delivered a rigorous definition of the greedy transfer operator \(\mathcal{L}_s\), proved its trace‑class nature for \(\operatorname{Re}s>1\), and outlined its analytic extension to the critical half‑plane.  We now have a bona fide nuclear operator whose Fredholm determinant can be studied.  The next programme will compute this determinant and relate it to the Riemann zeta function.

\begin{thebibliography}{9}
\bibitem{Mayer}
D.~H.~Mayer, \emph{The thermodynamic formalism approach to Selberg's zeta function for \(\mathrm{PSL}(2,\mathbb{Z})\)}, Bull. Amer. Math. Soc. (N.S.) 25 (1991), 55--60.
\bibitem{Efrat}
I.~Efrat, \emph{Dynamics of the continued fraction map and the spectral theory of \(\mathrm{SL}(2,\mathbb{Z})\)}, Invent. Math. 114 (1993), 207--218.
\bibitem{ProgrammeI}
V.~Geere et al., \emph{Harmonic Sine Transform Programme I}, 2026.
\bibitem{ProgrammeII}
V.~Geere et al., \emph{Harmonic Sine Transform Programme II}, 2026.
\bibitem{ProgrammeIII}
V.~Geere et al., \emph{Harmonic Sine Transform Programme III}, 2026.
\end{thebibliography}

\end{document}